%% foundations.tex — Mathesis: Foundations of a Tier Calculus for Scientific Claims
%%
%% Build:  latexmk -xelatex foundations.tex      (or ./build.sh)
%% Engine: XeLaTeX is REQUIRED. The source contains Unicode mathematics
%%         (ℚ ℝ α′ ω²r³ ‖∇u‖ ∀ ∃ ≤ →) passed through to unicode-math rather
%%         than transliterated into macros, so that the LaTeX source and the
%%         Lean source contain the SAME GLYPHS. That is the notation standard
%%         applied to itself (docs/NOTATION.md: one symbol, three columns that
%%         must never diverge). pdflatex will not compile this file.

\documentclass[11pt,a4paper]{article}

\usepackage{fontspec}
\usepackage{unicode-math}
\setmainfont{TeX Gyre Pagella}
\setmathfont{TeX Gyre Pagella Math}
\setmonofont{DejaVu Sans Mono}[Scale=0.82]

\usepackage[strings]{underscore}   % identifiers like Reff_ge_sqrt in text mode
\usepackage[a4paper,margin=2.6cm,bottom=3cm]{geometry}
\usepackage{booktabs}
\usepackage{longtable}
\usepackage{array}
\usepackage{listings}
\usepackage{xcolor}
\usepackage{microtype}
\usepackage{titlesec}
\usepackage{enumitem}
\usepackage{needspace}
\usepackage{fancyhdr}
\usepackage{mathesis}              % the programme notation standard — dogfooded
\usepackage[hidelinks,bookmarks=true]{hyperref}

\hypersetup{
  pdftitle={Mathesis: Foundations of a Tier Calculus for Scientific Claims},
  pdfauthor={SocrateAI Stream 0},
  pdfsubject={Epistemic bookkeeping for machine-checked science},
  colorlinks=true, linkcolor=mathesisink, urlcolor=mathesisink, citecolor=mathesisink,
}

\definecolor{mathesisink}{RGB}{20,60,110}
\definecolor{rulegray}{RGB}{120,120,120}
\definecolor{codebg}{RGB}{248,248,245}
\definecolor{leankw}{RGB}{120,20,120}
\definecolor{leancom}{RGB}{100,110,100}

%% ---------------------------------------------------------------------------
%% Lean listings
%% ---------------------------------------------------------------------------
\lstdefinelanguage{Lean}{
  morekeywords={theorem,def,lemma,example,axiom,sorry,by,have,show,exact,intro,
    structure,inductive,namespace,end,import,open,noncomputable,Prop,Type,
    trivial,unfold,simp,ring,field_simp,decide,norm_num,convert,induction,with,
    variable,instance,deriving,abbrev,protected,private},
  sensitive=true,
  morecomment=[l]{--},
  morecomment=[n]{/-}{-/},
  morestring=[b]",
}
\lstset{
  basicstyle=\ttfamily\footnotesize,
  keywordstyle=\color{leankw}\bfseries,
  commentstyle=\color{leancom}\itshape,
  backgroundcolor=\color{codebg},
  frame=leftline, framerule=1.1pt, rulecolor=\color{rulegray},
  framexleftmargin=6pt, xleftmargin=10pt,
  breaklines=true, breakatwhitespace=true,
  showstringspaces=false,
  columns=fullflexible,
  aboveskip=1.0em, belowskip=1.0em,
  literate=
    {ℚ}{{$\mathbb{Q}$}}1 {ℝ}{{$\mathbb{R}$}}1 {ℕ}{{$\mathbb{N}$}}1 {ℤ}{{$\mathbb{Z}$}}1
    {α}{{$\alpha$}}1 {β}{{$\beta$}}1 {γ}{{$\gamma$}}1 {δ}{{$\delta$}}1
    {ι}{{$\iota$}}1 {ω}{{$\omega$}}1 {τ}{{$\tau$}}1 {π}{{$\pi$}}1 {μ}{{$\mu$}}1
    {→}{{$\rightarrow$}}1 {≤}{{$\leq$}}1 {≥}{{$\geq$}}1 {≠}{{$\neq$}}1
    {∀}{{$\forall$}}1 {∃}{{$\exists$}}1 {∈}{{$\in$}}1 {¬}{{$\lnot$}}1
    {₀}{{$_0$}}1 {₁}{{$_1$}}1 {₂}{{$_2$}}1 {₃}{{$_3$}}1 {′}{{$'$}}1
    {√}{{$\sqrt{}$}}1 {∇}{{$\nabla$}}1 {∧}{{$\land$}}1 {≫}{{$\gg$}}1
}

%% ---------------------------------------------------------------------------
%% Section styling
%% ---------------------------------------------------------------------------
\titleformat{\section}{\Large\bfseries\color{mathesisink}}{\thesection}{0.7em}{}
\titleformat{\subsection}{\large\bfseries}{\thesubsection}{0.6em}{}
\titleformat{\subsubsection}{\normalsize\bfseries\itshape}{}{0em}{}
\titlespacing*{\section}{0pt}{2.2ex plus 1ex minus .2ex}{1.2ex plus .2ex}

\pagestyle{fancy}\fancyhf{}
\renewcommand{\headrulewidth}{0.2pt}
\fancyhead[L]{\footnotesize\color{rulegray}Mathesis — Foundations of a Tier Calculus}
\fancyhead[R]{\footnotesize\color{rulegray}\thepage}

\setlist[itemize]{leftmargin=1.3em,itemsep=0.25em,topsep=0.35em}
\setlist[enumerate]{leftmargin=1.5em,itemsep=0.25em,topsep=0.35em}

%% A boxed aside for the "what this does NOT claim" notes, which carry more
%% weight in this document than the claims do.
\newcommand{\caveat}[1]{%
  \par\vspace{0.6em}\noindent
  \begin{tabular}{@{}p{0.4em}@{\hspace{0.8em}}p{\dimexpr\linewidth-2.2em\relax}@{}}
  \textcolor{rulegray}{\rule{2pt}{\baselineskip}} & \small #1
  \end{tabular}\par\vspace{0.6em}}

\newcommand{\tierorder}{\ensuremath{\mathsf{X} < \mathsf{C} < \mathsf{L} < \mathsf{B} < \mathsf{A}}}

\begin{document}

\begin{titlepage}
\centering
\vspace*{2.5cm}
{\Huge\bfseries\color{mathesisink} Mathesis\par}
\vspace{0.8em}
{\LARGE Foundations of a Tier Calculus\\[0.3em] for Scientific Claims\par}
\vspace{2.2em}
{\large\itshape Stream 0 of SocrateAI\par}
\vspace{0.6em}
{\large Version 1.0 \quad·\quad 2026-08-14\par}
\vspace{3.2cm}

\begin{minipage}{0.82\textwidth}
\small
\textbf{Tier of this document: C} (governance and exposition).

It states no mathematical result of its own. Every technical claim it makes is
either a citation to a Tier~A artifact in this repository, or a Tier~C
observation about another repository, and is labelled as such.

\vspace{1em}

Seven scientific streams share mathematics and a verification technology. Before
this repository existed they shared no way of saying \emph{how strongly anything
is known}. Two of them used the letter \textsf{B} for incompatible things: one
meaning \emph{a program checked it}, the other meaning \emph{a referee checked
it}. This document sets out the notation, the kernel-verified theorem about it,
and the five worked cases that tested the apparatus before it was pointed at
anything that matters.
\end{minipage}

\vfill
{\footnotesize\color{rulegray}
\texttt{github.com/xaviercallens/SocrateAI-Scientific-Mathesis}\par
\vspace{0.4em}
All four verification gates green at the time of writing. Not signed off:\\
no human statement-adequacy audit has yet been performed.\par}
\end{titlepage}

\tableofcontents
\newpage

\section{Context}

SocrateAI runs seven concurrent scientific streams. They share mathematics — the T-dual
effective radius \texttt{Reff}, the Sym² recurrence lock, dyadic shell decompositions,
enstrophy functionals — and they share a verification technology, Lean~4. What they did
\emph{not} share, before this repository existed, was a way of saying \textbf{how strongly
anything is known}.

\begin{center}\small
\begin{tabular}{@{}clp{9.6cm}@{}}
\toprule
\# & Code & Object of study \\
\midrule
0 & \texttt{MX} & \textbf{Mathesis} — the notation, kernel, and ledger the others use \\
1 & \texttt{MF} & MechanicaFluidorum — 3D Navier–Stokes global regularity, via Hypothesis U \\
2 & \texttt{AE} & AutoEvolve — K3 selection in the DualScale K3×T² landscape \\
3 & \texttt{QK} & Quantum Agora — tensor-network, QUBO and Cirq execution of landscape search \\
4 & \texttt{HG} & Hypergraph — discrete hypergraph cosmology, computation-augmented generation \\
5 & \texttt{RM} & RajMath — the RAMA Ramanujan engine, mock modular forms, BPS entropy \\
6 & \texttt{TN} & TNN Univers Model — thermodynamic/topological neural networks \\
7 & \texttt{VD} & Videoo — scientific communication \\
\bottomrule
\end{tabular}
\end{center}

Stream~0 does no science. \textbf{Its object of study is the record.} The first theorem here is
about ledgers, not about geometry — and that is the point, because the record is the one
artifact every other stream produces and none of them owns.

\section{Motivation}

Not from first principles. From three failures found by looking.

\subsection{One letter, two meanings}

Stream~1's specification defines:

\begin{quote}\itshape
\textbf{B — Checkable}: identities validated in exact rational arithmetic; no floats.
\end{quote}

Stream~5's README defines:

\begin{quote}\itshape
\textbf{Tier B — Established}: Peer-reviewed literature, pinned to exact values.
\end{quote}

One means \emph{a program checked it}. The other means \emph{a referee checked it}. These are
not close. And artifacts already cross between those two repositories: Stream~1's Lean gate
compiles against Stream~5's working tree. The day a \textbf{claim} crosses with its letter
attached, a citation silently becomes a computation, and nothing in either repository notices.

This is \texttt{MX-C-0001}. It is the founding observation.

\subsection{A rule with no gate}

Stream~5's README states five CI rules, attributing them to \texttt{dualscale_ci.yml}:

\begin{quote}\itshape
\textbf{Rule R1}: No \texttt{axiom} declarations allowed in \texttt{.lean} files.
\end{quote}

That file does not exist. The repository has no \texttt{.github/} directory. A comment-stripped
scan of its Lean tree on 2026-08-13 found \textbf{34 \texttt{axiom} declarations across 14
files}, plus two \texttt{sorry}. The commit that introduced both \texttt{sorry}s is titled
\emph{``zero-sorry Lean 4 verification''}.

The sharpest instance, \texttt{DualScale/Physics/AubinLions.lean}, axiomatises the type, the
observable, the \textbf{conclusion predicate}, and the \textbf{implication} — then proves a
theorem by applying the axiom:

\begin{lstlisting}[language=Lean]
axiom VelocityField : Type                                   -- the type, opaque, no inhabitant
axiom enstrophy_of : VelocityField → ℝ                       -- the observable, opaque
axiom aubin_lions_compactness (S : Set VelocityField) : Prop -- the CONCLUSION, opaque
axiom aubin_lions_apply ... : aubin_lions_compactness S      -- the IMPLICATION, asserted

theorem aubin_lions_enstrophy_compactness ... : aubin_lions_compactness S := by
  exact aubin_lions_apply h_ens_bound h_ke_bound             -- the proof IS the axiom
\end{lstlisting}

Since \texttt{aubin_lions_compactness} has no definition, the theorem asserts nothing
falsifiable; and since \texttt{VelocityField} has no exhibited inhabitant, the hypotheses may be
vacuous. By that stream's \textbf{own Rule R3} — \emph{``logically vacuous propositions are
strictly prohibited and will fail the build''} — it fails. Nothing ran.

This is \texttt{MX-C-0003}. It is not an accusation. It is what happens to everyone when no gate
runs, and it is the most legible argument available for why Stream~0 exists.

\subsection{The same content, handled correctly, still failed}

Stream~1 met the \emph{same mathematics} — Aubin–Lions compactness — the stricter way: as named
hypothesis parameters, \textbf{never axioms}. It was audited anyway. The verdict, recorded in
that file's own header:

\begin{quote}\small
\textbf{TIER C (DRAFT), demoted 2026-08-13 by external audit.} The external audit's B1 verdict:
\texttt{AubinLionsStatement} and \texttt{ProdiSerrinStatement} as bare \texttt{Prop → Prop}
arrows \emph{``completely bypass PDE theory. The Lean kernel is merely verifying a logical
tautology (A → B → C).''} \textbf{Accepted.} This file remains kernel-compiled and gated (so it
cannot rot) but its claims are \textbf{TIER C} until repaired.
\end{quote}

A stream demoted its own flagship reduction and kept the file gated so the demotion could not be
quietly forgotten. That is the standard. Stream~0's job is to make it cheap and general.

\section{Narrative: Newton, Leibniz, and the century of ghosts}

The calculus was founded twice, independently, in the late seventeenth century. Both foundings
were correct. Both built on infinitesimal ideas reaching back through Cavalieri and Fermat to
Archimedes' method of exhaustion. The story of what happened next is the argument of this
document, and it has three acts.

\subsection*{Act I — Notation is not cosmetic}

Newton's \emph{method of fluxions} wrote the derivative of $x$ as $\dot{x}$ and reasoned about
quantities flowing in time. Leibniz wrote $dy/dx$ and $\int y\,dx$, and reasoned about ratios and
sums of differentials.

The two systems are mathematically equivalent. \textbf{Leibniz's notation won}, and it won for
reasons that have nothing to do with mathematical content:

\begin{itemize}
\item $dy/dx$ \emph{displays the chain rule}. $\frac{dy}{dt} = \frac{dy}{dx}\frac{dx}{dt}$ looks
like cancellation, so it is hard to get wrong. Fluxion notation gives no such scaffolding.
\item $\int y\,dx$ \emph{carries the variable of integration}. Substitution becomes bookkeeping
instead of insight.
\item Higher derivatives, partial derivatives and multiple integrals extend by pattern in the
Leibnizian system, and by invention in the Newtonian one.
\end{itemize}

British mathematics, loyal to Newton after the priority dispute, kept fluxions and fell behind
Continental analysis for roughly a century. It was not less able. It was worse equipped. The
Analytical Society at Cambridge in the 1810s campaigned explicitly for \emph{``the principles of
pure D-ism against the Dot-age of the University''} — a pun that was also a correct diagnosis.

\caveat{\textbf{The lesson Stream~0 takes:} the notation you record results in determines which
errors are easy to make and which are impossible to miss. \tierorder{} is a notation claim, not a
philosophy claim. It is chosen so that ``this Tier~A theorem cites a Tier~L result'' becomes
\emph{visibly} wrong — the way $dy/dx$ makes a botched chain rule visibly wrong.}

\subsection*{Act II — The dispute was an epistemics failure, and the adjudication was rigged}

The priority quarrel was not a disagreement about mathematics. Both men had it. It was a failure
of \emph{record-keeping}: no shared notation, no dated ledger of who established what and when,
no agreed standard of evidence for a claim of precedence.

In 1712 the Royal Society convened a committee to settle it and published the
\emph{Commercium Epistolicum}. The committee was appointed by the Society's president. Its report
was written, anonymously, by the same president — Newton. He then reviewed it, also anonymously.

An institution issued a verdict on a dispute in which its head was a party, and signed it with
its own authority. \textbf{The verdict happened to favour the person who was, on the evidence,
first to the method} — and it was still worthless as adjudication, because nothing independent
checked it.

This is exactly §2.2, three centuries earlier: an authority certifying itself, with no gate. It
is why agent self-reports are not evidence, why Gate~3 fails the build when two implementations
disagree \emph{without adjudicating which is right}, and why this repository is forbidden from
issuing scientific verdicts about the streams it serves.

\subsection*{Act III — A century of computing on foundations nobody had}

Newton and Leibniz both reasoned with infinitesimals that neither could justify. In 1734 George
Berkeley published \emph{The Analyst} and named the problem exactly:

\begin{quote}\itshape
And what are these Fluxions? The Velocities of evanescent Increments? And what are these same
evanescent Increments? They are neither finite Quantities, nor Quantities infinitely small, nor
yet nothing. May we not call them the Ghosts of departed Quantities?
\end{quote}

Berkeley was \textbf{right}. The objection stood, essentially unanswered, until Cauchy and
Weierstrass built the $\varepsilon$–$\delta$ foundations in the nineteenth century — roughly 150
years later. And Abraham Robinson's non-standard analysis (1961) eventually vindicated the
infinitesimals themselves, so the original intuition was sound all along.

Now: \textbf{during those 150 years, everyone kept computing, and it worked.} Euler produced an
enormous body of correct results on foundations that did not exist. Withholding the calculus
until it was rigorous would have been a catastrophe for science.

This is the subtlest and most important lesson, and it is the one a naive reading of a tier
system gets wrong. The tier calculus does \emph{not} say ``do not compute before you are
rigorous''. Euler at Tier~C was worth more than almost anyone at Tier~A. What it says is:

\begin{center}\large\bfseries Record which one you have.\end{center}

Tier~X is not a rebuke. It is a \emph{permission} — permission to explore with floats, sample,
guess and follow intuition, on the single condition that the result is labelled and cannot be
cited as though a machine had checked it. The failure mode is never ``someone conjectured''. It
is ``someone conjectured, and eighteen months later the conjecture was being cited as a theorem,
and nobody could point to when it changed status''.

Leibniz himself wanted the fix. His programme had two halves: the \emph{characteristica
universalis}, a universal formal notation for concepts; and the \emph{calculus ratiocinator}, a
mechanical procedure for deciding what follows from what. His famous promise for disputes —
\emph{``Let us calculate, Sir, and see who is right''} — \textbf{never arrived}. What arrived
instead, in the twentieth century, was the proof assistant: a \emph{calculus ratiocinator} that
actually runs. Stream~0 is the missing half, the \emph{characteristica}, built for the specific
job of recording scientific certainty rather than scientific content.

\section{Key concepts}

\subsection{The tier lattice}

Five tiers in one linear order:

\begin{center}\small
\begin{tabular}{@{}cll@{}}
\toprule
& Tier & Admission criterion \\
\midrule
\textsf{A} & Kernel      & Lean 4, zero \texttt{sorry}, axiom footprint matches its declared allowlist \\
\textsf{B} & Checkable   & finite statement decided in exact ℚ/ℤ, deterministic, shipping a \\
           &             & negative control \textbf{demonstrated to fail} \\
\textsf{L} & Literature  & peer-reviewed, cited to a \textbf{quoted theorem statement} — not an abstract \\
\textsf{C} & Conjecture  & proposal, analogy, unverified reduction, design memo \\
\textsf{X} & Exploratory & floats, sampling, plots, LLM output — \textbf{may never be cited} \\
\bottomrule
\end{tabular}
\end{center}

The order is about \textbf{how much of the checking a machine did} — not about mathematical
depth. A Tier~L theorem of Bourgain is deeper than every Tier~A row in this repository.

Tier~\textsf{L} is the tier that did not exist before. Its absence is what forced literature and
computation to share the letter \textsf{B} (§2.1), and its admission criterion — a \emph{quoted
theorem statement}, never an abstract or a retrieval summary — came from a real incident where a
tool-generated paraphrase conflicted with a stream's own description of its work.

\subsection{Evidence kinds cap tiers}

Orthogonal to the order, every claim declares \emph{how} it was established, and that caps what
it may claim:

\begin{center}\small
\begin{tabular}{@{}ll@{}}
\toprule
Evidence kind & Ceiling \\
\midrule
\texttt{lean_axioms}   & \textsf{A} \\
\texttt{exact_harness} & \textsf{B} \\
\texttt{citation}      & \textsf{L} \\
\texttt{argument}      & \textsf{C} \\
\texttt{numeric}, \texttt{llm_output} & \textsf{X} \\
\bottomrule
\end{tabular}
\end{center}

\texttt{llm_output} caps at \textsf{X} \textbf{by charter}. No model output promotes a tier —
including the output of the model that wrote this sentence.

\subsection{Soundness, and the theorem that matters}

A ledger assigns to each identifier a tier and a list of claims it rests on. It is \texttt{Sound}
when \emph{no claim is filed at a tier above any claim it \textbf{directly} cites}.

Every stream already checks that, informally, by eye. The kernel-verified content of Stream~0 is
that soundness on direct edges propagates through the \textbf{entire transitive closure}:

\begin{lstlisting}[language=Lean]
theorem tier_le_of_depends (hL : Sound L) :
    ∀ {a b : ι}, Depends L a b → (L a).tier ≤ (L b).tier

theorem no_kernel_claim_rests_on_weaker (hL : Sound L)
    (ha : (L a).tier = Tier.A) (hab : Depends L a b) : (L b).tier = Tier.A
\end{lstlisting}

\noindent(\texttt{MX-A-0003}, \texttt{MX-A-0004} — Mathlib-free, footprint \texttt{[]}.)

\textbf{Why this elementary result earns a kernel proof.} Transitivity is the property informal
bookkeeping loses \emph{first}. A chain \textsf{B} → \textsf{B} → \textsf{L} has its unsound
edge at the \emph{far} end — the middle claim cites literature it may not cite. A checker
validating each row against its own citations flags that middle row and says nothing about the
\textbf{head}, which is the row someone will actually cite. The transitive check is what reports
the head is contaminated too.

That is the retraction question — \emph{what else falls?} — and nobody answers it by eye, because
doing so requires holding the whole graph in mind, and the graph is exactly the thing that got too
big to hold.

The consequence people trip over: \textbf{a Tier~A claim may not cite a Tier~L theorem.} The
resolution is not to weaken the tier system — it is to take the literature result as an explicit
\textbf{hypothesis parameter}. The theorem becomes Tier~A \emph{conditionally}, the condition is
visible in its type, and the ledger records it.

\subsection{The model/world split}

This is the concept that makes the calculus usable in science rather than only in mathematics.

A scientific claim is almost always \textbf{two claims wearing one name}:

\begin{center}\small
\begin{tabular}{@{}clc@{}}
\toprule
 & Claim & Ceiling \\
\midrule
\textbf{M} & the mathematics follows from the model & \textsf{A} — provable \\
\textbf{W} & the model describes the system         & \textsf{C} at best — empirical \\
\bottomrule
\end{tabular}
\end{center}

``Hardy–Weinberg equilibrium'' names both the algebra and the assertion that a population is at
it. ``Conservation of energy'' names both a theorem about a Hamiltonian and a claim about an
apparatus. The Tier~A reputation of \textbf{M} silently underwrites \textbf{W}.

Under the calculus the split is structural rather than rhetorical: \textbf{W} is a separate row
that \emph{cites} \textbf{M}, and by \texttt{MX-A-0004} it can never be Tier~A — however good the
data is.

\subsection{Identifiers carry their tier}

\begin{center}
\texttt{<STREAM>-<TIER>-<NNNN>} \qquad e.g.\ \texttt{MF-A-0007}, \texttt{MX-C-0003}
\end{center}

The tier letter is \textbf{in the identifier}, so a promotion changes the identifier and every
stale citation elsewhere becomes \emph{lexically} wrong instead of \emph{silently} wrong. The old
row is kept and rewritten to point forward; \texttt{supersedes} records the trail.

\emph{Worked instance:} \texttt{MX-C-0004} (Lotka–Volterra bridge, OPEN) → \texttt{MX-A-0011}
(proved), 2026-08-13.

\section{Notation (the \emph{characteristica})}

One symbol table, three columns that must never diverge: LaTeX macro, meaning, Lean identifier.
This document is typeset with that package.

\textbf{First-use rule (normative):} a symbol may appear in prose only if its row exists in the
table.

\begin{center}\small
\begin{tabular}{@{}lll@{}}
\toprule
LaTeX & Meaning & Lean \\
\midrule
\verb|\Reff{\alpha'}{R}| & T-dual effective radius $\max(R,\alpha'/R)$ & \texttt{Mathesis.Scale.Reff} \\
\verb|\seam|             & fundamental length $\sqrt{\alpha'}$          & — derived \\
\verb|\HypU|             & Hypothesis U — \textbf{quantifier order normative} & \texttt{Mathesis.Statements.HypothesisU} \\
\verb|\tier{A}|          & epistemic tier tag                            & \texttt{schemas/ledger.schema.json} \\
\verb|\lean{Name}|       & the declaration a claim rests on              & — \\
\bottomrule
\end{tabular}
\end{center}

\textbf{Reserved vocabulary.} These words carry gate-backed meanings:

\begin{center}\small
\begin{tabular}{@{}lp{9.5cm}@{}}
\toprule
Word & Reserved for \\
\midrule
\textbf{theorem}     & Tier~A only \\
\textbf{verified}    & a kernel or exact-arithmetic gate accepted it, here, reproducibly \\
\textbf{established} & Tier~A or Tier~L — never a Tier~B instance check \\
\textbf{proven}      & Tier~A. Not ``we checked 10\,000 cases'' \\
\textbf{uniform}     & never without its quantifier order stated explicitly \\
\bottomrule
\end{tabular}
\end{center}

The package enforces this structurally: the \texttt{mtheorem} environment takes a
\textbf{mandatory} Lean declaration name, and \texttt{mclaim} takes a \textbf{mandatory} tier. An
untiered claim in a programme paper is exactly what it exists to prevent.

\caveat{\textbf{On quantifier order.} Hypothesis U is $\exists C\,\forall\text{cutoff}$, not
$\forall\text{cutoff}\,\exists C$. Its entire content is uniformity in the cutoff. Swap the
quantifiers and every symbol still typechecks, every proof still compiles, and the statement has
become a much weaker one. No kernel catches that. This is why quantifier order is normative in
prose and why the human audit exists.}

\section{Solver principles}

\subsection{What ``solver'' means here}

\begin{center}\itshape
A solver is a producer of claims that a dumb exact checker can replay.
\end{center}

Not a thing that computes an answer — a thing that emits a record another, simpler program can
independently re-decide. A certificate that cannot be replayed by a checker simpler than the
producer is not a certificate.

\subsection{The four gates}

\begin{center}\small
\begin{tabular}{@{}clp{6.4cm}@{}}
\toprule
Gate & Checks & Demonstrated to fail \\
\midrule
1 & every \texttt{tier_b_*.py} exits 0 & by construction — the controls are its probes \\
2 & Lean compiles; no \texttt{sorry}; footprints as declared & planted \texttt{sorry}; planted indented axiom \\
3 & Python and Rust checkers agree on 30 corpus cases & caught a real divergence on its first run \\
4 & ledger \texttt{Sound}; both ledgers name the same ids & planted orphan id \\
\bottomrule
\end{tabular}
\end{center}

\textbf{A gate that has only ever been observed passing is indistinguishable from a gate that
cannot fail.} Every one of the four has been shown to reject a planted defect, and each
demonstration is recorded.

\subsection{The footprint is the gate}

Never accept ``there is no \texttt{sorry} in the source'' as evidence. A failed proof still
defines the theorem name in the environment; the source looks clean and only
\texttt{\#print axioms} reveals the \texttt{sorryAx}. Stream~1 caught two broken proofs this way
that had been \emph{reported as passing}.

\subsection{Two implementations, no shared dependency}

The ledger checker exists in Python and in Rust, each written against the specification rather
than against the other. The Rust crate has \textbf{zero} crates.io dependencies, including for
JSON — a differential gate whose two sides call the same library tests that library once and the
ledger logic zero times. When they disagree, \textbf{neither is trusted} until a human
adjudicates.

\subsection{The human audit is not automatable}

\textbf{A theorem is Tier~A the moment the kernel accepts it. It is \emph{citable} only after a
human has audited the statement.} The kernel confirms the proof establishes the statement;
nothing mechanical confirms the statement is the one anyone meant.

This is the only step in the pipeline whose output cannot be predicted from the pipeline's own
signals — which is exactly why it cannot be removed to go faster.

\subsection{No dependency Stream 0 cannot rebuild in a minute}

The Lean core is Mathlib-free and cold-builds in about two seconds with no \texttt{import} at
all. The Rust crate has no dependencies. The Python package is standard-library only. Stream~0's
gate is what every other stream will eventually call; it must never be the reason another stream
is blocked.

\section{Methodology: the loop}

\begin{center}\ttfamily\small
0 SPLIT $\rightarrow$ 1 REFUTE $\rightarrow$ 2 CHECK $\rightarrow$ 3 PROVE $\rightarrow$
4 \textbf{HUMAN} $\rightarrow$ 5 FILE
\end{center}

\textbf{0 — Split.} Write \textbf{M} and \textbf{W} as separate sentences. Only \textbf{M} enters
the Lean file.

\textbf{1 — Refute before proving.} Perturb a coefficient, zero a variable, try the degenerate
case, in exact arithmetic at small parameters. A false conjecture absorbs unbounded proof effort
and returns nothing; a five-second sweep kills most. Stream~1 found a convolution formula,
sourced from a web search, that was \textbf{identically zero} — caught at $M = 1,2,3$, and
invisible to any amount of proof effort because the statement was perfectly provable and
perfectly useless.

\textbf{2 — Check.} Exhaustive over a small exact grid beats random over a large one: a failure
names a reproducible tuple rather than a seed. Ship a control targeting the \textbf{plausible
near-miss}.

\textbf{3 — Prove.} Declare the gate contract in the header. Work over ℚ when the harness uses
\texttt{Fraction}, so the theorem and the check are about \emph{one} structure rather than two
that agree numerically.

\textbf{4 — Audit.} Human. Three questions: does the formal statement say what the gloss claims;
are the hypotheses non-vacuous; \textbf{are the quantifiers where the science needs them}.

\textbf{5 — File.} Artifact, both ledger rows, and gate output, in one commit.

\caveat{\textbf{Health is the kill ratio, not the promotion count.} A good run refutes most of
what it proposes. A loop that promotes everything is broken, and a loop that has never produced a
rejection has not been tested.}

\section{Domains}

The apparatus is domain-agnostic because it constrains \emph{evidence}, not \emph{content}. What
varies by domain is where the model/world seam falls and which tier is reachable at all.

\begin{center}\small
\begin{tabular}{@{}lccp{5.4cm}@{}}
\toprule
Domain & \textbf{M} ceiling & \textbf{W} ceiling & Characteristic trap \\
\midrule
Pure mathematics       & \textsf{A} & — & vacuous statements; junk values ($x/0 = 0$) \\
Theoretical physics    & \textsf{A}$^*$ & \textsf{L} & fixing a quantified parameter as a global constant \\
Computational physics  & \textsf{B} & \textsf{C} & floats certifying a rounding mode, not a law \\
Population biology     & \textsf{A} & \textsf{C} & the model's assumptions hold of nothing real \\
Genomics               & \textsf{B} & \textsf{C} & statistical significance read as mechanism \\
Machine learning       & \textsf{X} & \textsf{X} & a validated model treated as evidence about the world \\
Quantum hardware       & \textsf{X} & \textsf{X} & sampled, noisy, non-deterministic by construction \\
\bottomrule
\end{tabular}
\end{center}
\noindent{\footnotesize $^*$ conditional on the force or action law, entered as a hypothesis.}

Two rows deserve comment.

\textbf{Machine learning and quantum hardware cap at \textsf{X}, and this is not a slight.} It is
what lets them be cited \emph{for what they actually establish}. A quantum annealer's output can
\textbf{steer} a landscape search; it can never \textbf{support} a claim. A trained model
reproducing a dynamic is Tier~X evidence about the model, which is exactly what it is.

\textbf{Computational physics is where the float ban does its work.} ``Agrees to $10^{-15}$''
certifies a rounding mode. The discipline is to reformulate until the statement is exact — and
§9's Kepler case shows that reformulating usually \emph{improves} the statement rather than
restricting it.

\section{The five use cases}

Five claims, three sciences, increasing complexity, all Tier~A over ℚ or ℝ with an
exact-arithmetic harness checking the same statements. They are deliberately elementary: the
point is that the \emph{apparatus} survives contact with mathematics nobody will argue about, so
that when it is pointed at something contested, a failure is a failure of the claim and not of
the plumbing.

\begin{center}\small
\begin{tabular}{@{}clllp{4.3cm}@{}}
\toprule
\# & Domain & Claim & Artifact & Exercises \\
\midrule
1 & Mathematics & $1+3+\dots+(2n-1) = n^2$ & \texttt{OddSums} & the pipeline itself \\
2 & Physics & elastic collision conserves $p$ \emph{and} $E$ & \texttt{ElasticCollision} & exact vs.\ ``agrees to $10^{-15}$'' \\
3 & Biology & Hardy–Weinberg allele invariance & \texttt{HardyWeinberg} & \textbf{model vs.\ world} \\
4 & Physics & Kepler III from inverse-square & \texttt{Kepler} & the float ban improving a statement \\
5 & Biology & Lotka–Volterra conservation & \texttt{LotkaVolterra\{,Flow\}} & \textbf{deferral, then promotion} \\
\bottomrule
\end{tabular}
\end{center}

\subsection{UC1 — Sum of odd numbers}

Mathlib-free, 0.8\,s. Chosen because a failure can only be the pipeline.

\emph{Taught:} the declared footprint went stale within minutes. The header said
\texttt{allow=} and the proof needed \texttt{omega}, which carries
\texttt{[propext, Quot.sound]}. A proof rewrite was attempted to reach \texttt{[]}, failed, and
the \textbf{declaration} was amended — the direction the rule requires.

\subsection{UC2 — Elastic collision}

Both momentum and kinetic energy conserved exactly, over 1296 enumerated mass/velocity
combinations, for $m_1 + m_2 \neq 0$.

\emph{The control that matters:} a perfectly \textbf{inelastic} collision conserves momentum but
\textbf{not} energy. A momentum-only implementation passes it. That control is what makes the
energy check mean something.

\emph{Witness:} at $m_2 = -m_1$ the formulae divide by zero, Lean's junk-value convention returns
$0$, and momentum is \textbf{not} conserved — which is why the side condition is in both
statements.

\subsection{UC3 — Hardy–Weinberg: the model/world split}

Allele frequency is invariant under random mating; equilibrium is reached in \textbf{one}
generation from arbitrary starting genotype frequencies.

The algebra is Tier~A. \emph{``This population is at Hardy–Weinberg equilibrium''} is Tier~C at
best — the model assumes infinite population, random mating, no selection, mutation, migration or
drift, and not one holds of any real population. In most write-ups one phrase covers both, and
the Tier~A reputation of the first underwrites the second.

\emph{The control that matters:} selection against the recessive homozygote \textbf{moves} the
allele frequency. Hardy–Weinberg exists to make selection visible; a check that could not detect
a violation would make the null model unfalsifiable.

\subsection{UC4 — Kepler III: the float ban earning its keep}

$T^2 \propto a^3$ carries $\pi$, which is irrational, so nothing is exactly checkable and the
standard response is floats. But \textbf{$\pi$ was never carrying any physics.} With the reduced
period $\tau = T/2\pi = 1/\omega$:
\[
  \frac{r^3}{\tau^2} \;=\; \omega^2 r^3 \;=\; GM
\]
Exactly rational whenever $r$ and $GM$ are; the constant becomes $GM$ rather than $4\pi^2/GM$;
$\pi$ has left the statement. Nothing approximated, nothing lost — $T^2 \propto a^3$ follows
because $T$ and $\tau$ differ by the same factor for every orbit.

\caveat{\textbf{The constraint did not make the physics harder to state. It located a constant
that was decoration and removed it.}}

\emph{The control that matters:} under an inverse-\textbf{cube} force, $\omega^2 r^3$ stops being
constant — so the check is about gravity, not arithmetic.

\subsection{UC5 — Lotka–Volterra: deferral, then promotion}

\emph{First half.} The algebraic core — substituting the vector field into $dV/dt$ and clearing
the logarithmic derivatives leaves an expression identically zero. Four terms cancelling in
pairs. Verified over 15\,625 exact parameter combinations.

The analytic step — that this expression \textbf{is} $dV/dt$ — was filed \textbf{OPEN at Tier~C},
refusing both shortcuts: \texttt{axiom} (forbidden), and the conditional restatement
\emph{``given $dV/dt = D$, and $D = 0$, therefore $dV/dt = 0$''} — which compiles, carries no
axioms, looks like Tier~A, and is a tautology of exactly the shape Stream~1 was demoted for.

\emph{Second half.} \textbf{The deferral was wrong, and finding out why was the more useful
result.} The calculus modules were absent from one Mathlib checkout and present in the other.
Pointed at the second, \texttt{conserved_along_flow} compiles on the first attempt, in nine
lines. $V$ is now proved genuinely constant along any trajectory: derivative exactly $0$, not
``$0$ given that it is $0$''.

\texttt{MX-C-0004} → \texttt{MX-A-0011}. New identifier, \texttt{supersedes} recorded, old row
retained pointing forward. \textbf{This is the first tier promotion in the ledger} — until it
happened, the promotion rule was schema with no instance.

\caveat{\emph{The refusals were right regardless.} Had the calculus genuinely been missing,
\texttt{axiom} and the tautological conditional would both still have been wrong answers.}

\section{Benefit, measured (not asserted)}

The honest measure of an apparatus is not how many claims it certifies but \textbf{whether it has
recently caught something its author believed}. By that measure:

\subsection{Defects found in the gates, by the campaign}

\begin{enumerate}
\item \textbf{Gate 2 was silently skipping theorems.} Lean wraps \texttt{\#print axioms} output
past its line width; the parser was line-anchored and dropped every wrapped declaration.
\texttt{ElasticCollision.lean} reported \emph{``1 footprint''} for a file with two theorems,
\textbf{and passed}. Because wrapping is triggered by long names, the bug hid precisely the
deeply-namespaced declarations in new modules. It is a false \emph{all-clear}, in the gate that
is the sole evidence for every Tier~A row. Fixed; the gate now cross-checks
\texttt{\#print axioms} directives against footprints parsed, and runs parser self-tests on every
invocation.

\item \textbf{Two theorems had never been footprint-checked} — \texttt{Tier.X_le} and
\texttt{witnessChain_reach} had no \texttt{\#print axioms} line since they were written.

\item \textbf{The axiom scanner was evadable by one space.} It matched \texttt{\^{}axiom} at
column 0, so an indented, \texttt{private}, or attributed axiom all passed. Its ``ignores prose''
control passed \emph{for the same reason it was evadable}. One property was bought with the other
and the suite could not tell.

\item \textbf{Three non-vacuity witnesses were wrong on first attempt} — parameters chosen by
inspection landed on degenerate cases, and one ``counterexample'' was simply false ($\tfrac12 +
\tfrac12$ does sum to 1). All caught by the kernel with $\vdash$ \texttt{False}.

\item \textbf{A deferral rested on an unchecked premise.} A missing Lean module reports as an
unknown identifier — indistinguishable from a nonexistent lemma or a false statement. An unbuilt
module masqueraded as a mathematical obstruction.
\end{enumerate}

Items 1, 2 and 3 would have stayed invisible without the campaign. That is the argument for
running elementary cases through new apparatus \textbf{before} pointing it at anything that
matters.

\subsection{Findings about the streams}

The Tier~B collision; the CI file that does not exist while 34 axioms and two \texttt{sorry} sit
in the tree it claims to gate; and two Mathlib subsets that differ, neither complete, so that
which one is resolved determines what is \emph{provable}.

A fourth, verified directly on 2026-08-13, is the same shape in two more streams:
\texttt{TNN/specs/TNN_Invariants.lean:46} contains \texttt{sorry} in
\texttt{energy_conservation}, while that stream's dashboard hardcodes
\texttt{"lean4_status": "PROVEN (Zero-Sorry)"} \textbf{fifteen times as a string literal},
computed from nothing. And \texttt{RajMathRecovery/NAMAGIRI.lean} proves Hypothesis U and
singularity prevention \texttt{by trivial} over \texttt{Prop := True}, with \texttt{Real} defined
as \texttt{Float}.

\caveat{\textbf{The recurring pattern across four of seven streams is not bad mathematics. It is
a status string that no check produces.} A README asserting a CI file that does not exist; a
dashboard field that is a literal; a docstring describing what a \texttt{trivial} proof would
establish if its predicate were not \texttt{True}. In every case the \emph{code} is honest and
the \emph{label} is not, and there is no gate between them. That gap is the entire product.}

\subsection{One finding came from \emph{not} trusting a scanner}

A survey subagent reported ``839 \texttt{sorry} and 264 \texttt{axiom}'' in Stream~5's tree.
Re-measured excluding the vendored Mathlib: \textbf{2 and 34}. The agent had counted Mathlib
itself, where \texttt{sorry} appears in test files and where \texttt{axiom} is how
\texttt{propext} and \texttt{Classical.choice} are \emph{defined}. The figure was specific,
confident, and made the case being argued stronger — three reasons to check it harder, all of
which push the other way.

\subsection{What is deliberately not claimed}

That any of this is true of the world. Every Tier~A row is a statement about a model whose
assumptions are written into the theorem's type. \textbf{A green board says the bookkeeping
holds. It licenses nothing about collisions, populations, planets, fluids, or the universe.}

\section{Future work: the hard problems}

This section is \textbf{Tier~C throughout}. It describes what the apparatus would contribute to
five hard problems. It does not claim progress on any of them, and Stream~0 is forbidden from
issuing scientific verdicts about the streams it serves.

\caveat{Hard problems fail at the \textbf{seams} — where an unproven analytic step is quietly
absorbed, where a quantifier is fixed that should have stayed free, where a Tier~C model
assumption inherits the authority of the Tier~A algebra built on top of it. Mathesis does not
attack the problem. It makes the seams \textbf{visible, typed, and mechanically checked}, so that
effort is spent where the difficulty actually is.}

\subsection{Navier–Stokes, and how to avoid the singularity}

Global regularity for 3D Navier–Stokes: do smooth initial data remain smooth for all time, or can
enstrophy blow up in finite time? Stream~1 attacks it via \textbf{Hypothesis U} — uniform-in-$\alpha'$
enstrophy control of the frequency-truncated system:
\[
  \sup_{\alpha' > 0}\; \sup_{0 \le t \le T}\; \bigl\| \nabla u^{(\alpha')}(t) \bigr\|_{L^2(\mathbb{T}^3)} \;<\; \infty
\]

\textbf{The quantifier order is the content.} $\exists C\,\forall\alpha'$ says one constant works
for every truncation — which is what lets $\alpha' \to 0$ recover the true equation.
$\forall\alpha'\,\exists C$ says each truncation has its own constant, which is nearly free and
implies nothing. \textbf{Both typecheck. Both compile. No kernel distinguishes them.}

Stream~5's \texttt{TDuality.lean} shows the failure mode live: it declares
\texttt{noncomputable def alpha_prime : ℝ := 1}. With $\alpha'$ pinned, every statement about the
$\alpha' \to 0$ limit is unstateable and nothing reports it.

\subsubsection{Singularity avoidance — the mechanism, already Tier A}

The dual-scale programme's geometric idea is that the effective radius seen by the truncated
system is $\Reff{\alpha'}{R} = \max(R, \alpha'/R)$, which \textbf{cannot go below $\seam$}. Below
the seam the scale \emph{bounces}. There is no path to zero scale. Stream~0 has proved this,
unconditionally and axiom-free (\texttt{MX-A-0005}):

\begin{center}\small
\begin{tabular}{@{}lp{9.2cm}@{}}
\toprule
Theorem & Content \\
\midrule
\texttt{Reff_pos}                 & the effective radius is never zero \\
\texttt{Reff_ge_sqrt}             & $\sqrt{\alpha'} \le \Reff{\alpha'}{R}$ — \textbf{a universal floor} \\
\texttt{Reff_eq_sqrt_iff}         & the floor is attained at \textbf{exactly one} point, the self-dual radius \\
\texttt{Reff_bounce}              & below the seam, $\Reff{\alpha'}{R} = \alpha'/R$ — the reflection \\
\texttt{genesis_no_singularity}   & $0 < \texttt{tDualRadius}\ \alpha'\ R$ — no collapse \\
\texttt{one_div_sq_le_of_sqrt_le} & $\sqrt{\alpha'} \le x \Rightarrow 1/x^2 \le 1/\alpha'$ \\
\bottomrule
\end{tabular}
\end{center}

That last one is the bridge to the fluid estimates: a minimum length becomes a \textbf{maximum}
on the $1/x^2$ factors that drive an enstrophy cascade.

\subsubsection{The failure mode, in the wild}

\texttt{RajMathRecovery/NAMAGIRI.lean} — a standalone file at that repository's root — contains:

\begin{lstlisting}[language=Lean]
def Real := Float
def uniformBoundedness (D : EnstrophyFunctional) : Prop := True
/-- Ramanujan's sum of tails algebraically bounds the ultra-high-frequency
    modes, proving Hypothesis U. -/
theorem hypothesis_U_bound (D : EnstrophyFunctional) : uniformBoundedness D := by
  trivial

def topologicalFracture (state : Complex) : Prop := True
theorem prevents_singularity (res : Nat) : topologicalFracture (fractionLimit res) := by
  trivial
\end{lstlisting}

The predicates are \texttt{True}. The proofs are \texttt{trivial}. \texttt{Real} is
\texttt{Float}. The docstrings claim Hypothesis U and singularity prevention — the two central
objects of the Navier–Stokes programme.

\caveat{\textbf{In fairness, and precisely:} the file is imported by nothing and appears in no
lakefile, so it contaminates no build and no gated claim rests on it. It is a skeleton, and
writing skeletons is legitimate. What is missing is the \emph{label}. Under the calculus this is
Tier~X — \texttt{Prop := True} is the definition of a vacuous statement, and Stream~5's own Rule
R3 forbids exactly it. The gap between a sketch and a claim is one docstring wide, and nothing
currently sits in it.}

\subsection{Quantum wave, T-duality, K3×T² and F-theory}

Select the K3×T² geometry reproducing observed cosmological parameters — treating the landscape
as a search space rather than an anthropic accident.

This is the domain most exposed to \textbf{tier inflation}, because candidates are produced
\emph{mechanically and in volume}. A search result is Tier~X until something checks it, and the
gap between ``the search proposed it'' and ``we verified it'' is where a landscape programme
loses its footing.

\textbf{Quantum hardware output is Tier~X, mechanically.} A D-Wave anneal or a Cirq run can
\textbf{steer} a search; it can never \textbf{support} a claim. This is not a limitation imposed
on Stream~3 — it is what lets its results be cited for what they establish.

\textbf{The unresolved question, stated rather than assumed.} Is a search-produced candidate that
passes an exact-arithmetic filter \emph{Tier~B}, or \emph{Tier~X with a Tier~B filter-result
attached}? These are different claims. Stream~0 has not decided it; it is reserved to the owner.

\subsection{Black hole entropy}

Reproduce the Bekenstein–Hawking entropy by microstate counting. Stream~5's route is BPS state
counting through Ramanujan-type asymptotics.

The seam falls between the \textbf{counting function} and the \textbf{entropy}. The asymptotic
growth of a partition-like generating function is a hard but ordinary analytic fact. That this
growth \emph{is} the entropy of a black hole is a physical identification resting on a chain of
dualities. These are \textbf{M} and \textbf{W}, routinely written as one sentence.

The entropy function itself is concrete and reasonable —
\texttt{macroscopicEntropy (c_eff n) := 2 * π * √(c_eff * n / 6)}, the Cardy formula. It is
the \emph{chain around it} that is axiomatised: \texttt{bps_macroscopic_entropy},
\texttt{rademacher_first_order}, \texttt{saddle_point_asymptotic},
\texttt{bps_implies_spectral_gap}, \texttt{hagedorn_divergence}, \texttt{modular_s_transform},
\texttt{rogers_ramanujan_identity}.

That last is a \textbf{theorem} — proved, famous, and available in the literature. As an axiom it
contributes nothing but risk: Tier~L material entered at a level that permanently pollutes every
downstream footprint.

\textbf{The fix is mechanical and loses nothing.} Convert each \texttt{axiom} to a hypothesis
parameter. The theorems survive, become explicitly conditional, and the dependence becomes
visible in the type instead of invisible in the footprint. Results with genuine literature
backing then move to Tier~\textsf{L} — the tier that did not exist before, and precisely the one
this material needs.

\subsection{Dark matter and dark energy}

The most evidence-rich and model-poor domain in physics: a large, consistent body of observation
and a wide space of models fitting it.

Nowhere else is the model/world gap wider or more consequential. The \textbf{observations} are
Tier~L — peer-reviewed, pinned to published values with quoted uncertainties. The \textbf{fits}
are Tier~B at best, and only if the fitting is exact and deterministic, which it usually is not.
The \textbf{interpretations} are Tier~C, and they are what gets reported.

The characteristic failure is not fabrication. It is a Tier~C model selection acquiring, over a
few years of citation, the felt authority of the Tier~L observations it was fitted to.

Model comparison becomes a ledger query: competing models citing the same Tier~L observational
rows, each at its own tier, make \emph{``what does this rest on that the alternative does
not?''} answerable by traversing \texttt{supports} rather than by literature review. And a
simulation is Tier~X — an N-body run is sampled and floating-point. It may steer; it may not
support.

\caveat{\textbf{Honest limit.} Stream~4's tree has \textbf{not} been audited. What is known is
only its shape: it carries \texttt{dark_matter/}, \texttt{proofs/}, \texttt{lean_oracle/}, an
MFDM continuum-limit paper and NANOGrav SGWB specifications — and it is the only one of the seven
that is \textbf{not a git repository}, so it has no commit history to audit against. Its tier
vocabulary, Lean inventory and axiom counts are \textbf{unmeasured}. The absence of a finding is
not a clean bill.}

\subsection{DNA and genetics}

From sequence to mechanism: which variants cause which phenotypes, and by what pathway. No
SocrateAI stream works on this \emph{yet}, though two have written it into their forward plans:
Stream~6's frontier names \emph{``scRNA-seq topological networks \dots multi-gene epigenetic
regulation \dots in oncology''}, and Stream~5 has a design note proposing to \emph{``treat DNA
folding or the propagation of genetic mutations as a cellular automaton''}.

UC3 is already the entry point. The Hardy–Weinberg algebra is Tier~A; \emph{``this population is
at Hardy–Weinberg equilibrium''} is Tier~C. Every genome-wide association study rests on that
distinction, and the characteristic failure is a \textbf{statistical association read as a
mechanism} — a Tier~B correlation cited as a Tier~A causal claim.

\texttt{llm_output} caps at \textsf{X}, and this domain needs it most: sequence models are now
routinely used to predict variant effects. A model's confidence is Tier~X evidence \emph{about
the model}. It is not evidence about a genome, however well the model validates.

Much of classical genetics is Tier~B material currently reported in floats — segregation ratios,
linkage distances, pedigree probabilities and Hardy–Weinberg expectations are all exact rational
arithmetic.

\section{Conclusion}

Leibniz's programme had two halves. The \emph{calculus ratiocinator} — a mechanical procedure for
deciding what follows from what — arrived three centuries late, as the proof assistant, and it
works. The \emph{characteristica universalis} — the notation those procedures operate on — did
not arrive, and its absence is what this repository addresses.

Not a universal notation for \emph{concepts}. Leibniz's ambition there was too large and probably
incoherent. A universal notation for \textbf{certainty}: five tiers, one linear order, an
orthogonal evidence kind that caps what each row may claim, and one kernel-verified theorem
saying that soundness on direct citations propagates through the entire transitive closure.

The theorem is elementary. That is deliberate. It earns a kernel proof not because it is hard but
because \textbf{transitivity is the property informal bookkeeping loses first} — every stream
checks direct citations by eye, and none checks the closure, because the closure is exactly the
thing that grew too big to hold in mind.

Three things are worth carrying away.

\textbf{First, the notation is the mechanism.} Leibniz's $dy/dx$ beat Newton's $\dot{x}$ not by
being more correct but by making a botched chain rule \emph{visibly} wrong, and a century of
British mathematics paid for the difference. \tierorder{} is chosen on the same principle: to
make ``this Tier~A theorem cites a Tier~L result'' a lexical error rather than a judgement call.

\textbf{Second, low tiers are permissions, not rebukes.} Newton and Leibniz computed for 150 years
on foundations Berkeley correctly demolished, and they were right to. Euler at Tier~C was worth
more than almost anyone at Tier~A. The tier calculus does not ask anyone to wait for rigour. It
asks them to \textbf{record which one they have} — because the failure was never that someone
conjectured, but that eighteen months later the conjecture was being cited as a theorem and
nobody could say when it changed status.

\textbf{Third, the apparatus must be tested on things nobody disputes.} Five elementary use cases
across mathematics, physics and biology found five defects — three in the gates themselves,
including a false all-clear in the gate that is the sole evidence for every Tier~A row, and one
deferral resting on a premise nobody had checked. Three of those would have stayed invisible.
Pointing an unvalidated instrument at a Millennium Problem tells you nothing about the problem.

What Stream~0 refuses is as load-bearing as what it provides. It issues no scientific verdicts.
It does not automate the human statement-adequacy audit — the one step whose result cannot be
predicted from the pipeline's own signals, which is precisely why it cannot be removed to go
faster. It does not adjudicate when two independent implementations disagree; it fails the build
and escalates to a person. And it does not claim that a green board licenses anything beyond the
internal consistency of a record.

The 1712 \emph{Commercium Epistolicum} reached the right conclusion about priority and was
worthless as adjudication, because the president of the deciding body wrote the verdict on his
own dispute and signed it with the institution's authority. \textbf{Being right is not the same
as being checkable.} A programme that cannot tell the difference will eventually be unable to
tell whether it is right.

\vspace{1em}
\begin{center}\itshape
Calculemus was never a promise that the calculation would succeed —\\
only that when it did not, everyone would be able to see exactly where.
\end{center}

\appendix
\section{Artifacts}

\begin{center}\small
\begin{tabular}{@{}lp{8.4cm}@{}}
\toprule
Kernel core        & \texttt{lean/Mathesis/TierCalculus.lean} — Mathlib-free, footprint \texttt{[]} \\
Consolidated Reff  & \texttt{lean/Mathesis/Scale/Reff.lean} \\
Use cases          & \texttt{lean/Mathesis/Applications/} — 5 cases, 6 modules \\
Reference checker  & \texttt{python/mathesis/} — standard library only, no floats \\
Independent checker& \texttt{rust/mathesis-verify/} — zero dependencies, including JSON \\
Gates              & \texttt{scripts/verify.sh} — four gates, each demonstrated to fail \\
Ledger             & \texttt{LEDGER.md} + \texttt{ledger.jsonl} — 23 rows, \texttt{Sound} \\
Lessons            & \texttt{LL.md} — 15 entries, each with its incident \\
Notation           & \texttt{latex/mathesis.sty}, \texttt{docs/NOTATION.md} \\
Owner decisions    & \texttt{docs/OWNER_BRIEF.md} \\
\bottomrule
\end{tabular}
\end{center}

\vspace{0.6em}
\noindent\textbf{Status.} All four gates green. \textbf{Not signed off}: no human
statement-adequacy audit has been performed, so Stream~0's own ledger currently carries the very
gap it exists to close in others, and says so.

\end{document}
